\section{Trace norm}

\begin{definition}[Schatten one-norm]\label{def:schatten_one_norm}
    \lean{Matrix.schattenOneNorm}
    \leanok
    For a matrix $A\in\MN{D}$, the singular values $s_0(A),s_1(A),\ldots$ form
    a finitely supported family: only finitely many are nonzero. The Schatten
    one-norm of $A$ is their sum, that is, the sum of the finitely many nonzero
    singular values:
    \begin{align}
        \lVert A\rVert_1
         & =\sum_{i\,:\,s_i(A)\neq 0}s_i(A).
        \label{eq:entropy_schatten_one_norm}
    \end{align}
    This is the $p=1$ case of the Schatten $p$-norm
    \cite[Chapter~8, Section~8.1]{Wolf2012Quantum}. The equivalent closed
    formula $\lVert A\rVert_1=\sum_{i=0}^{D-1}s_i(A)$, summing over all $D$
    singular values including trailing zeros, is part of
    Theorem~\ref{thm:trace_norm_elementary}.
\end{definition}

\begin{definition}[Trace norm]\label{def:trace_norm}
    \lean{Matrix.traceNorm}
    \leanok
    \uses{def:schatten_one_norm}
    The trace norm of $A\in\MN{D}$ is its Schatten one-norm:
    \begin{align}
        \lVert A\rVert_{\tr}
         & =\lVert A\rVert_1.
        \label{eq:entropy_trace_norm}
    \end{align}
\end{definition}

Wolf records the equivalent formula
$\lVert A\rVert_1=\tr\lvert A\rvert$
\cite[Chapter~8, Section~8.1]{Wolf2012Quantum}; this is
Theorem~\ref{thm:trace_norm_abs}.

\begin{theorem}[Finite singular-value expansions]\label{thm:trace_norm_singular_values}
    \lean{Matrix.schattenOneNorm_eq_sum_support}
    \lean{Matrix.traceNorm_eq_sum_support}
    \lean{Matrix.traceNorm_eq_sum_range_finrank_range}
    \leanok
    \uses{def:schatten_one_norm, def:trace_norm}
    The Schatten one-norm and trace norm are the sums over the finite support of
    the singular-value sequence. Equivalently, the trace norm is the sum over
    the indices below the rank of the represented linear map.
\end{theorem}

\begin{proof}\leanok
    The singular-value sequence satisfies $s_i(A)=0$ for all
    $i\geq\rank(A)$, so summing over the support and summing
    over $\{0,\ldots,\rank(A)-1\}$ give the same value.
\end{proof}

\begin{theorem}[Elementary trace-norm properties]
    \label{thm:trace_norm_elementary}
    \lean{Matrix.traceNorm_eq_sum_fin}
    \lean{Matrix.schattenOneNorm_nonneg}
    \lean{Matrix.traceNorm_nonneg}
    \lean{Matrix.schattenOneNorm_eq_zero_iff}
    \lean{Matrix.traceNorm_eq_zero_iff}
    \lean{Matrix.traceNorm_zero}
    \lean{Matrix.traceNorm_pos_iff}
    \leanok
    \uses{def:trace_norm, def:schatten_one_norm, thm:trace_norm_singular_values}
    For every $A\in\MN{D}$,
    \begin{align}
        \lVert A\rVert_{\tr}
         & =\sum_{i=0}^{D-1}s_i(A),
        \label{eq:entropy_trace_norm_fin} \\
        \lVert A\rVert_{\tr}
         & \geq 0.
        \notag
    \end{align}
    Also, $\lVert A\rVert_{\tr}=0$ if and only if $A=0$.
    Equivalently, $\lVert A\rVert_{\tr}>0$ if and only if $A\neq0$.
\end{theorem}

\begin{proof}\leanok
    The singular values satisfy $s_i(A)\geq0$, so
    $\lVert A\rVert_{\tr}=0
        \iff \forall i,\,s_i(A)=0
        \iff A=0$.
\end{proof}

\begin{theorem}[Trace norm as trace of the absolute value]
    \label{thm:trace_norm_abs}
    \lean{Matrix.traceNorm_eq_re_trace_abs,
        Matrix.rectangularTraceNorm_eq_re_trace_abs}
    \leanok
    \uses{def:trace_norm}
    For every $A\in\MN{D}$, with
    $\lvert A\rvert=\sqrt{A^\dagger A}$ the positive-semidefinite square root
    of $A^\dagger A$,
    \begin{align}
        \lVert A\rVert_{\tr}
         & =\re(\tr\lvert A\rvert).
        \label{eq:entropy_trace_norm_abs}
    \end{align}
    This is the formula $\lVert A\rVert_1=\tr[\lvert A\rvert]$
    of \cite[Chapter~8, Section~8.1]{Wolf2012Quantum}; the trace of the
    positive-semidefinite matrix $\lvert A\rvert$ is real.
\end{theorem}

\begin{proof}\leanok
    \uses{lem:trace_norm_sqrt_eigenvalues}
    The map $T^\dagger T$, for $T$ the linear map on $\C^D$ represented by
    $A$, is represented by $A^\dagger A$, so the singular values of $A$ are
    $s_i(A)=\sqrt{\lambda_i}$ with $\lambda_0,\ldots,\lambda_{D-1}$ the
    eigenvalues of $A^\dagger A$. Diagonalizing
    $A^\dagger A=U\diag(\lambda_i)U^\dagger$ gives
    $\lvert A\rvert=U\diag(\sqrt{\lambda_i})U^\dagger$, hence
    \begin{align}
        \lVert A\rVert_{\tr}
         & =\sum_{i=0}^{D-1}s_i(A)
        =\sum_{i=0}^{D-1}\sqrt{\lambda_i}
        =\tr\lvert A\rvert
        =\re(\tr\lvert A\rvert),
        \notag
    \end{align}
    where the last step uses that $\lvert A\rvert$ is positive semidefinite,
    so $\tr\lvert A\rvert\geq0$ is real.
\end{proof}

\begin{theorem}[Trace-norm homogeneity]
    \label{thm:trace_norm_homogeneous}
    \lean{Matrix.traceNorm_smul}
    \leanok
    \uses{def:trace_norm}
    For every $c\in\C$ and $A\in\MN{D}$,
    $\lVert cA\rVert_{\tr}
        =\lvert c\rvert\,\lVert A\rVert_{\tr}$.
    This is the homogeneity axiom for matrix norms
    \cite[Chapter~8, Section~8.1]{Wolf2012Quantum}.
\end{theorem}

\begin{proof}\leanok
    \uses{thm:trace_norm_abs}
    From $(cA)^\dagger(cA)=\lvert c\rvert^2\,A^\dagger A$ and uniqueness of
    the positive-semidefinite square root,
    $\lvert cA\rvert=\lvert c\rvert\,\lvert A\rvert$. Linearity of the trace
    gives
    $\tr\lvert cA\rvert
        =\lvert c\rvert\,\tr\lvert A\rvert$, and the claim follows
    from Theorem~\ref{thm:trace_norm_abs}.
\end{proof}

\begin{theorem}[Trace-norm unitary invariance]
    \label{thm:trace_norm_unitary_invariance}
    \lean{Matrix.traceNorm_unitary_mul}
    \lean{Matrix.traceNorm_mul_unitary}
    \lean{Matrix.traceNorm_unitary_mul_unitary}
    \leanok
    \uses{def:trace_norm}
    For all unitaries $U,V\in\MN{D}$ and every $A\in\MN{D}$,
    \begin{align}
        \lVert UA\rVert_{\tr}
         & =\lVert AV\rVert_{\tr}
        =\lVert UAV\rVert_{\tr}
        =\lVert A\rVert_{\tr}.
        \notag
    \end{align}
    The trace norm is thus unitarily invariant
    \cite[Chapter~8, Section~8.1]{Wolf2012Quantum}.
\end{theorem}

\begin{proof}\leanok
    \uses{thm:trace_norm_abs}
    From $(UA)^\dagger(UA)=A^\dagger U^\dagger UA=A^\dagger A$ the absolute
    values agree: $\lvert UA\rvert=\lvert A\rvert$. For the right factor,
    $(AV)^\dagger(AV)=V^\dagger(A^\dagger A)V$, and since
    $V^\dagger\lvert A\rvert V$ is positive semidefinite with
    \begin{align}
        (V^\dagger\lvert A\rvert V)^2
         & =V^\dagger\lvert A\rvert\,VV^\dagger\,\lvert A\rvert V
        =V^\dagger(A^\dagger A)V,
        \notag
    \end{align}
    uniqueness of the positive-semidefinite square root gives
    $\lvert AV\rvert=V^\dagger\lvert A\rvert V$. Cyclicity of the trace then
    yields
    $\tr\lvert AV\rvert
        =\tr(\lvert A\rvert\,VV^\dagger)
        =\tr\lvert A\rvert$. For the two-sided case,
    $\lVert UAV\rVert_{\tr}
        =\lVert UA\rVert_{\tr}
        =\lVert A\rVert_{\tr}$
    by applying the right and left cases in turn.
\end{proof}

\begin{lemma}[Trace norm from the spectrum of $A^\dagger A$]
    \label{lem:trace_norm_sqrt_eigenvalues}
    \lean{Matrix.traceNorm_eq_sum_sqrt_eigenvalues_adjoint_comp_self}
    \leanok
    \uses{def:trace_norm}
    For every $A\in\MN{D}$, the trace norm is the sum of the square roots of
    the eigenvalues of the positive operator $A^\dagger A$:
    \begin{align}
        \lVert A\rVert_{\tr}
         & =\sum_{i=0}^{D-1}\sqrt{\lambda_i(A^\dagger A)}.
        \notag
    \end{align}
\end{lemma}

\begin{proof}\leanok
    \uses{thm:trace_norm_elementary}
    Since $s_i(A)=\sqrt{\lambda_i(A^\dagger A)}$ by the definition of singular
    value, the finite singular-value expansion
    (\ref{eq:entropy_trace_norm_fin}) gives
    \begin{align}
        \lVert A\rVert_{\tr}
         & =\sum_{i=0}^{D-1}s_i(A)
        =\sum_{i=0}^{D-1}\sqrt{\lambda_i(A^\dagger A)}.
        \notag
    \end{align}
\end{proof}

\begin{lemma}[Trace norm from the matrix eigenvalues of $A^\dagger A$]
    \label{lem:trace_norm_sqrt_matrix_eigenvalues}
    \lean{Matrix.traceNorm_eq_sum_sqrt_eigenvalues}
    \leanok
    \uses{def:trace_norm}
    For every $A\in\MN{D}$, with $\lambda_0,\ldots,\lambda_{D-1}$ the
    eigenvalues of the matrix $A^\dagger A$,
    $\lVert A\rVert_{\tr}
        =\sum_{i=0}^{D-1}\sqrt{\lambda_i}$.
\end{lemma}

\begin{proof}\leanok
    \uses{thm:trace_norm_abs}
    Diagonalizing $A^\dagger A=V\diag(\lambda_i)V^\dagger$ gives
    $\lvert A\rvert=V\diag(\sqrt{\lambda_i})V^\dagger$, so
    $\tr\lvert A\rvert=\sum_{i=0}^{D-1}\sqrt{\lambda_i}$,
    and the claim follows from Theorem~\ref{thm:trace_norm_abs}.
\end{proof}

\begin{theorem}[Trace norm and Hilbert--Schmidt norm comparison]
    \label{thm:trace_norm_frobenius_comparison}
    \lean{Matrix.frobenius_norm_le_traceNorm,
        Matrix.traceNorm_le_sqrt_dim_mul_frobenius_norm}
    \leanok
    \uses{lem:trace_norm_sqrt_matrix_eigenvalues}
    For every $A\in\MN{D}$,
    \begin{align}
        \lVert A\rVert_2
         & \leq \lVert A\rVert_1
        \leq \sqrt D\,\lVert A\rVert_2.
        \label{eq:trace_norm_frobenius_comparison}
    \end{align}
    The first inequality is the $p=1$, $p'=2$ case of
    \cite[Eq.~(8.1)]{Wolf2012Quantum}; the second is the corresponding case
    of \cite[Eq.~(8.7)]{Wolf2012Quantum}. Both are used in the proof of
    trace-norm convergence toward asymptotic states.
\end{theorem}

\begin{proof}\leanok
    Write $s_0,\ldots,s_{D-1}\geq0$ for the singular values of $A$.
    Then $\lVert A\rVert_2^2=\sum_i s_i^2$ and
    $\lVert A\rVert_1=\sum_i s_i$. The first inequality follows by expanding
    $(\sum_i s_i)^2$ and using nonnegativity. The second is the
    Cauchy--Schwarz bound
    $(\sum_i s_i)^2\leq D\sum_i s_i^2$.
\end{proof}

\begin{theorem}[Variational unitary formula for the trace norm]
    \label{thm:trace_norm_variational}
    \lean{Matrix.norm_trace_conjTranspose_mul_unitary_le}
    \lean{Matrix.exists_mem_unitaryGroup_trace_conjTranspose_mul_eq}
    \lean{Matrix.isGreatest_norm_trace_conjTranspose_mul_unitary}
    \lean{Matrix.traceNorm_eq_sSup_norm_trace_conjTranspose_mul_unitary}
    \leanok
    \uses{def:trace_norm}
    For every $A\in\MN{D}$,
    \begin{align}
        \lVert A\rVert_{\tr}
         & =\max\left\{\bigl|\tr[A^\dagger U]\bigr|
        \middle|  UU^\dagger=\Id\right\},
        \label{eq:entropy_trace_norm_variational}
    \end{align}
    and the maximum is attained: some unitary $U$ satisfies
    $\tr[A^\dagger U]=\lVert A\rVert_{\tr}$.
    See \cite[Chapter~8, Eq.~(8.11)]{Wolf2012Quantum}.
\end{theorem}

\begin{proof}\leanok
    \uses{lem:trace_norm_sqrt_matrix_eigenvalues, thm:unitary_polar_correction}
    For the upper bound, let $v_0,\ldots,v_{D-1}$ be an orthonormal
    eigenbasis of $A^\dagger A$ with eigenvalues $\lambda_i$, and let $U$ be
    unitary. Expanding the trace in this basis and applying the
    Cauchy--Schwarz inequality,
    \begin{align}
        \bigl|\tr[A^\dagger U]\bigr|
         & =\Bigl|\sum_i\langle Av_i,\,Uv_i\rangle\Bigr|
        \leq\sum_i\lVert Av_i\rVert\,\lVert Uv_i\rVert
        =\sum_i\sqrt{\lambda_i}
        =\lVert A\rVert_{\tr},
        \notag
    \end{align}
    since $\lVert Av_i\rVert^2=\langle v_i,\,A^\dagger A\,v_i\rangle
        =\lambda_i$ and $\lVert Uv_i\rVert=1$.

    For attainment, the polar decomposition
    (Theorem~\ref{thm:unitary_polar_correction}) writes
    $A=U\sqrt{A^\dagger A}$ with $U$ unitary. Since $\sqrt{A^\dagger A}$ is
    self-adjoint, $A^\dagger U=\sqrt{A^\dagger A}\,U^\dagger U
        =\sqrt{A^\dagger A}$, and
    $\tr\sqrt{A^\dagger A}=\sum_i\sqrt{\lambda_i}=\lVert A\rVert_{\tr}$.
\end{proof}

\begin{theorem}[Trace-norm duality over any finite index]
    \label{thm:trace_norm_duality_any_index}
    \lean{Matrix.exists_mem_unitaryGroup_trace_conjTranspose_mul_eq_trace_sqrt,
        Matrix.re_trace_sqrt_conjTranspose_mul_self_le_of_forall_norm_trace_mul_le,
        Matrix.l2_opNorm_unitary_le_one,
        Matrix.traceNorm_eq_re_trace_sqrt,
        Matrix.trace_sqrt_conjTranspose_mul_self_eq_traceNorm}
    \leanok
    \uses{def:trace_norm}
    Let $A$ be a square complex matrix indexed by any finite set. Some unitary
    $U$ satisfies $\tr[A^\dagger U]=\tr\sqrt{A^\dagger A}$. Consequently, if
    $\bigl|\tr(OA)\bigr|\leq\epsilon$ for every matrix $O$ of operator
    norm at most one, then $\tr\sqrt{A^\dagger A}\leq\epsilon$. For
    $A\in\MN{D}$, $\tr\sqrt{A^\dagger A}=\lVert A\rVert_{\tr}$.
\end{theorem}

\begin{proof}\leanok
    \uses{thm:unitary_polar_correction}
    Write $A=U\sqrt{A^\dagger A}$ with $U$ unitary
    (Theorem~\ref{thm:unitary_polar_correction}). Then
    $A^\dagger U=\sqrt{A^\dagger A}$ and $U^\dagger A=\sqrt{A^\dagger A}$.
    A unitary matrix has operator norm at most one, so $O=U^\dagger$ gives
    $\tr\sqrt{A^\dagger A}=\tr(U^\dagger A)\leq\epsilon$.
\end{proof}

\begin{theorem}[Trace-norm triangle inequality]
    \label{thm:trace_norm_triangle}
    \lean{Matrix.traceNorm_add_le}
    \leanok
    \uses{def:trace_norm}
    For all $A,B\in\MN{D}$,
    $\lVert A+B\rVert_{\tr}
        \leq\lVert A\rVert_{\tr}+\lVert B\rVert_{\tr}$.
    Together with homogeneity (Theorem~\ref{thm:trace_norm_homogeneous}) and
    definiteness (Theorem~\ref{thm:trace_norm_elementary}), this completes
    the norm axioms of \cite[Chapter~8, Section~8.1]{Wolf2012Quantum} for the
    trace norm.
\end{theorem}

\begin{proof}\leanok
    \uses{thm:trace_norm_variational}
    Choose by Theorem~\ref{thm:trace_norm_variational} a unitary $U$ with
    $\tr[(A+B)^\dagger U]=\lVert A+B\rVert_{\tr}$.
    Splitting the trace and applying the upper-bound half of the same
    theorem to $A$ and to $B$ separately,
    \begin{align}
        \lVert A+B\rVert_{\tr}
         & =\bigl|\tr[A^\dagger U]
        +\tr[B^\dagger U]\bigr|
        \leq\bigl|\tr[A^\dagger U]\bigr|
        +\bigl|\tr[B^\dagger U]\bigr|
        \leq\lVert A\rVert_{\tr}+\lVert B\rVert_{\tr}.
        \notag
    \end{align}
\end{proof}

\begin{lemma}[Jordan trace-norm formula]
    \label{lem:trace_norm_jordan}
    \lean{Matrix.IsHermitian.traceNorm_eq_re_trace_posPart_add_negPart}
    \leanok
    \uses{def:trace_norm}
    Let $H\in\MN{D}$ be Hermitian with Jordan decomposition $H=H^+-H^-$
    into orthogonal positive parts, $H^\pm\geq0$ and $H^+H^-=0$. Then
    $\lVert H\rVert_{\tr}=\tr[H^+]+\tr[H^-]$.
\end{lemma}

\begin{proof}\leanok
    \uses{thm:trace_norm_abs}
    The sum $H^++H^-$ is positive semidefinite, and by orthogonality its
    square is
    \begin{align}
        (H^++H^-)^2
         & =(H^+)^2+(H^-)^2
        =(H^+-H^-)^2
        =H^\dagger H,
        \notag
    \end{align}
    so $H^++H^-=\sqrt{H^\dagger H}=\lvert H\rvert$. Hence
    $\lVert H\rVert_{\tr}
        =\tr\lvert H\rvert
        =\tr[H^+]+\tr[H^-]$
    by Theorem~\ref{thm:trace_norm_abs}.
\end{proof}

\begin{lemma}[Support projection of the positive part]
    \label{lem:positive_part_support_projection}
    \lean{Matrix.IsHermitian.posPartSupportProj}
    \lean{Matrix.IsHermitian.posPartSupportProj_posSemidef}
    \lean{Matrix.IsHermitian.posPartSupportProj_le_one}
    \lean{Matrix.IsHermitian.posPartSupportProj_mul_posPartSupportProj}
    \lean{Matrix.IsHermitian.posPartSupportProj_mul_self}
    \leanok
    Let $A\in\MN{D}$ be Hermitian with positive part $A^+$. There is a
    matrix $\Pi$ with $0\leq\Pi\leq\Id$, $\Pi^2=\Pi$, and $\Pi A=A^+$, namely
    the orthogonal projection onto the support space of $A^+$.
\end{lemma}

\begin{proof}\leanok
    Diagonalize
    $A=U\diag(\lambda_1,\ldots,\lambda_D)\,U^\dagger$ and set
    \begin{align}
        \Pi
         & =U\diag
        (\chi(\lambda_1),\ldots,\chi(\lambda_D))\,U^\dagger,
        \notag
    \end{align}
    where $\chi$ is the indicator function of $(0,\infty)$. Each claim is
    read off eigenvalue-wise: $0\leq\chi\leq1$ gives
    $0\leq\Pi\leq\Id$, $\chi^2=\chi$ gives $\Pi^2=\Pi$, and
    $\chi(\lambda) \lambda=\max(\lambda,0)$ gives $\Pi A=A^+$.
\end{proof}

\begin{lemma}[Positive-semidefinite trace pairing bounds]
    \label{lem:psd_trace_pairing_bounds}
    \lean{Matrix.PosSemidef.re_trace_mul_nonneg}
    \lean{Matrix.PosSemidef.re_trace_mul_le_of_le_one}
    \leanok
    Let $X,C\in\MN{D}$ with $X\geq0$. If $C\geq0$, then
    $\tr[CX]\geq0$; if $C\leq\Id$, then
    $\tr[CX]\leq\tr[X]$.
\end{lemma}

\begin{proof}\leanok
    \uses{lem:psd_trace_product_nonneg}
    The first bound is Lemma~\ref{lem:psd_trace_product_nonneg}. For the
    second, $\tr[X]-\tr[CX]=\tr[(\Id-C)X]\geq0$
    by the same lemma, since $\Id-C\geq0$.
\end{proof}

\begin{lemma}[Positive maps preserve hermiticity]
    \label{lem:positive_map_hermitian}
    \lean{Matrix.isHermitian_map_of_positive}
    \leanok
    \uses{def:positive_map_rectangular}
    If $T:\MN{D}\to\MN{D'}$ is a positive linear map and $H\in\MN{D}$ is
    Hermitian, then $T(H)$ is Hermitian.
\end{lemma}

\begin{proof}\leanok
    \uses{def:positive_map_rectangular}
    Write $H=H^+-H^-$ with $H^\pm\geq0$. Then $T(H)=T(H^+)-T(H^-)$ is a
    difference of positive semidefinite matrices, hence Hermitian.
\end{proof}

\begin{lemma}[Common projection estimate]
    \label{lem:posPart_support_projection_estimate}
    \lean{Matrix.re_trace_posPart_map_le_aux}
    \leanok
    \uses{def:positive_map_rectangular}
    Let $T:\MN{D}\to\MN{D'}$ be a positive linear map and let
    $H\in\MN{D}$ be Hermitian.  Then
    $\tr[(TH)^+]\leq\tr[T(H^+)]$.
\end{lemma}

\begin{proof}\leanok
    \uses{lem:positive_map_hermitian,
        lem:psd_trace_pairing_bounds,
        lem:positive_part_support_projection}
    Let $\Pi_+$ be the support projection of $(TH)^+$ from
    Lemma~\ref{lem:positive_part_support_projection}, so that
    $\Pi_+ T(H) = (T(H))^+$ and $0\leq\Pi_+\leq\Id$.  The estimate follows
    from the chain
    \begin{align}
        \tr[(TH)^+] & = \tr[\Pi_+ T(H^+-H^-)] \notag                 \\
                    & = \tr[\Pi_+ T(H^+)] - \tr[\Pi_+ T(H^-)] \notag \\
                    & \leq \tr[\Pi_+ T(H^+)] \notag                  \\
                    & \leq \tr[T(H^+)] ,
    \end{align}
    using Lemmas~\ref{lem:psd_trace_pairing_bounds}
    and~\ref{lem:positive_map_hermitian}.
\end{proof}

\begin{lemma}[Positive-part trace estimate]
    \label{lem:positive_part_trace_estimate}
    \lean{Matrix.re_trace_posPart_map_le}
    \leanok
    \uses{def:positive_map_rectangular, def:trace_preserving_rectangular}
    Let $T:\MN{D}\to\MN{D'}$ be a trace-preserving positive linear map and
    let $H\in\MN{D}$ be Hermitian, with Jordan decompositions $H=P_+-P_-$
    and $T(H)=Q_+-Q_-$. Then $\tr[Q_+]\leq\tr[P_+]$.
\end{lemma}

\begin{proof}\leanok
    \uses{lem:posPart_support_projection_estimate,
        def:trace_preserving_rectangular}
    Lemma~\ref{lem:posPart_support_projection_estimate} gives
    $\tr[(TH)^+]\leq\tr[T(H^+)]$.  The result follows because
    trace preservation gives $\tr[T(H^+)]=\tr[H^+]$.
\end{proof}

\begin{theorem}[Trace-norm contractivity]
    \label{thm:trace_norm_contractivity}
    \lean{Matrix.traceNorm_map_le_of_positive_of_tracePreserving}
    \leanok
    \uses{def:trace_norm, def:positive_map_rectangular,
        def:trace_preserving_rectangular}
    Let $T:\MN{D}\to\MN{D'}$ be a trace-preserving positive linear map.
    Then for all Hermitian $H\in\MN{D}$,
    $\lVert T(H)\rVert_{\tr}\leq\lVert H\rVert_{\tr}$.
    See \cite[Chapter~8, Theorem~8.16]{Wolf2012Quantum}.
\end{theorem}

\begin{proof}\leanok
    \uses{lem:trace_norm_jordan, lem:positive_part_trace_estimate,
        lem:positive_map_hermitian}
    Write $H=P_+-P_-$ and $T(H)=Q_+-Q_-$ for the Jordan decompositions.
    Applying Lemma~\ref{lem:positive_part_trace_estimate} to $H$ and to
    $-H$ gives $\tr[Q_+]\leq\tr[P_+]$ and
    $\tr[Q_-]\leq\tr[P_-]$, so by the Jordan
    trace-norm formula (Lemma~\ref{lem:trace_norm_jordan}),
    \begin{align}
        \lVert T(H)\rVert_{\tr}
         & =\tr[Q_+]+\tr[Q_-]
        \leq\tr[P_+]+\tr[P_-]
        =\lVert H\rVert_{\tr}.
        \notag
    \end{align}
\end{proof}

\begin{corollary}[Trace-norm contractivity on states]
    \label{cor:trace_norm_contractivity_states}
    \lean{Matrix.traceNorm_map_sub_map_le_of_positive_of_tracePreserving}
    \leanok
    \uses{def:trace_norm, def:positive_map_rectangular,
        def:trace_preserving_rectangular, def:density_matrices}
    Let $T:\MN{D}\to\MN{D'}$ be a trace-preserving positive linear map.
    Then for all density matrices $\rho_1,\rho_2\in\MN{D}$,
    $\lVert T(\rho_1)-T(\rho_2)\rVert_{\tr}
        \leq\lVert\rho_1-\rho_2\rVert_{\tr}$.
    See \cite[Chapter~8, Eq.~(8.80)]{Wolf2012Quantum}.
\end{corollary}

\begin{proof}\leanok
    \uses{thm:trace_norm_contractivity}
    By linearity $T(\rho_1)-T(\rho_2)=T(\rho_1-\rho_2)$, and
    $\rho_1-\rho_2$ is Hermitian as a difference of positive semidefinite
    matrices, so Theorem~\ref{thm:trace_norm_contractivity} applies.
\end{proof}

\begin{lemma}[Trace-norm contraction coefficient]
    \label{lem:trace_norm_contraction_coefficient}
    \lean{Matrix.traceNorm_contraction_coefficient_eq_half_sSup_orthogonal,
        Matrix.orthogonalPureStateTraceNorms_nonempty_of_density}
    \leanok
    \uses{def:trace_norm, def:density_matrices}
    Let $T:\MN{D}\to\MN{D'}$ be a complex-linear map. Then
    \begin{align}
        \sup_{\substack{\rho_1,\rho_2\in\mc{D}_D                    \\ \rho_1\neq\rho_2}}
        \frac{\lVert T(\rho_1)-T(\rho_2)\rVert_{\tr}}
        {\lVert\rho_1-\rho_2\rVert_{\tr}}
        =\frac12\sup_{\substack{\lVert\psi\rVert=\lVert\phi\rVert=1 \\
                \langle\psi,\phi\rangle=0}}
        \left\lVert T\left(\ketbra{\psi}{\psi}
        -\ketbra{\phi}{\phi}\right)\right\rVert_{\tr}.
        \label{eq:trace_norm_contraction_coefficient}
    \end{align}
    The left supremum is over distinct density matrices in $\MN{D}$, and the
    right supremum is over orthogonal unit vectors in $\C^D$. No positivity or
    trace-preservation assumption is imposed on $T$.
    This is \cite[Chapter~8, Lemma~8.3, Eq.~(8.81)]{Wolf2012Quantum}.
\end{lemma}

\begin{proof}\leanok
    \uses{lem:trace_norm_jordan, thm:trace_norm_homogeneous,
        thm:trace_norm_triangle}
    An orthogonal pair of pure states gives a pair of distinct density matrices
    at trace-norm distance $2$, which proves the lower bound. For the reverse
    inequality, set $H=\rho_1-\rho_2$ and write its Jordan decomposition as
    $H=H^+-H^-$. Since $H$ is nonzero and traceless,
    $\tr[H^+]=\tr[H^-]=t>0$. Thus $P=H^+/t$ and $Q=H^-/t$ are density matrices
    with orthogonal supports, and homogeneity reduces the quotient for $H$ to
    $\frac12\lVert T(P-Q)\rVert_{\tr}$.

    Take spectral decompositions
    $P=\sum_i\lambda_i\ketbra{\psi_i}{\psi_i}$ and
    $Q=\sum_j\mu_j\ketbra{\phi_j}{\phi_j}$. Orthogonality of the supports gives
    $\psi_i\perp\phi_j$ whenever $\lambda_i\mu_j\neq0$, while
    $\lambda_i,\mu_j\geq0$ and $\sum_i\lambda_i=\sum_j\mu_j=1$. The
    product-weight expansion
    \begin{align}
        P-Q=\sum_{i,j}\lambda_i\mu_j
        \left(\ketbra{\psi_i}{\psi_i}-\ketbra{\phi_j}{\phi_j}\right)
        \notag
    \end{align}
    is therefore a convex combination of differences of orthogonal pure
    states. The triangle inequality and homogeneity bound its image under $T$
    by the right-hand supremum. Taking the supremum over
    $\rho_1\neq\rho_2$ proves the equality.
\end{proof}

\begin{definition}[Trace-norm distance to the asymptotic image]
    \label{def:trace_norm_asymptotic_distance}
    \lean{Matrix.traceNormAsymptoticDistance}
    \leanok
    \uses{def:trace_norm, def:peripheral_spectral_projections}
    For $T:\MN{D}\to\MN{D}$, let $T_\phi$ be its peripheral spectral
    projection and define
    \begin{align}
        \Delta_T(\rho):=\lVert\rho-T_\phi(\rho)\rVert_{\tr}.
    \end{align}
    This is Wolf's notation in
    \cite[Chapter~8, Proposition ``Convergence towards asymptotic
        states'', Eq.~(8.112)]{Wolf2012Quantum}.
\end{definition}

\begin{lemma}[Asymptotic-distance numerator identities]
    \label{lem:trace_norm_asymptotic_distance_numerator}
    \lean{Matrix.peripheralProjection_comp_pow,
        Matrix.peripheralWeightedProjection_apply_sub_peripheralProjection,
        Matrix.peripheralWeightedProjection_pow_apply_sub_peripheralProjection,
        Matrix.traceNormAsymptoticDistance_pow_eq,
        Matrix.traceNormAsymptoticDistance_pow_eq_sub_peripheralWeightedProjection_pow}
    \leanok
    \uses{def:trace_norm_asymptotic_distance,
        thm:peripheral_spectral_projection_properties}
    Let $T_\varphi:=T\circ T_\phi=T_\phi\circ T$ be Wolf's asymptotic
    dynamics. For every matrix $\rho$, every $n\geq0$, and every $n>0$ in
    the second equality,
    \begin{align}
        \Delta_T(T^n(\rho))
         & =\lVert T^n(\rho-T_\phi(\rho))\rVert_{\tr}               \\
         & =\lVert(T^n-T_\varphi^n)(\rho-T_\phi(\rho))\rVert_{\tr}.
        \label{eq:trace_norm_asymptotic_distance_numerator}
    \end{align}
    These are the numerator identities used in
    \cite[Chapter~8, Eq.~(8.114)]{Wolf2012Quantum}. They follow from
    $T_\phi T=TT_\phi=T_\varphi=T_\varphi T_\phi$; the positive-iterate
    condition records that the zeroth power of $T_\varphi$ is the identity.
\end{lemma}

\begin{proof}\leanok
    The commutation $T_\phi T=TT_\phi$ iterates to
    $T_\phi T^n=T^nT_\phi$, which proves the first equality by linearity.
    Idempotence of $T_\phi$ gives
    $T_\varphi(\rho-T_\phi(\rho))=0$; every positive power of
    $T_\varphi$ therefore vanishes on the same remainder, proving the second
    equality.
\end{proof}

\begin{lemma}[Pointwise form of the contraction-coefficient bound]
    \label{lem:trace_norm_contraction_pointwise}
    \lean{Matrix.traceNorm_map_sub_div_traceNorm_le_half_sSup_orthogonal}
    \leanok
    \uses{def:trace_norm, def:density_matrices}
    Let $S:\MN{D}\to\MN{D'}$ be complex-linear and let
    $\rho_1\neq\rho_2$ be density matrices. Then
    \begin{align}
        \frac{\lVert S(\rho_1)-S(\rho_2)\rVert_{\tr}}
        {\lVert\rho_1-\rho_2\rVert_{\tr}}
        \leq\frac12
        \sup_{\substack{\lVert\psi\rVert=\lVert\phi\rVert=1 \\
                \psi\perp\phi}}
        \left\lVert S\left(\ketbra{\psi}{\psi}
        -\ketbra{\phi}{\phi}\right)\right\rVert_{\tr}.
    \end{align}
    This is the pointwise inequality from Wolf Lemma~8.3 used in
    Equation~(8.115).
\end{lemma}

\begin{proof}\leanok
    \uses{lem:trace_norm_contraction_coefficient}
    The proof is the upper-bound half of
    Lemma~\ref{lem:trace_norm_contraction_coefficient}: normalize the positive
    and negative parts of $\rho_1-\rho_2$ to orthogonally supported density
    matrices and expand both in eigenprojectors.
\end{proof}

\begin{lemma}[Hilbert--Schmidt estimates for the asymptotic bound]
    \label{lem:asymptotic_hilbert_schmidt_estimates}
    \lean{Matrix.frobenius_norm_pureStateProj_sub_eq_sqrt_two,
        Matrix.sSup_orthogonalPureStateTraceNorms_le_hilbertSchmidtOperatorNorm}
    \leanok
    \uses{def:trace_norm,
        thm:trace_norm_frobenius_comparison}
    For $A\in\MN{D}$ and a complex-linear map $S:\MN{D}\to\MN{D}$,
    \begin{align}
        \lVert A\rVert_{\tr} & \leq\sqrt D\,\lVert A\rVert_2,              \\
        \lVert S(A)\rVert_2  & \leq\lVert S\rVert_{2\to2}\lVert A\rVert_2.
    \end{align}
    If $\psi$ and $\phi$ are orthogonal unit vectors, then
    $\lVert\ketbra{\psi}{\psi}-\ketbra{\phi}{\phi}\rVert_2=\sqrt2$.
    Consequently the orthogonal-pure-state supremum for $S$ is at most
    $\sqrt D\,\lVert S\rVert_{2\to2}\sqrt2$. These are precisely the
    estimates in Wolf Equations~(8.115)--(8.116).
\end{lemma}

\begin{proof}\leanok
    The first inequality is finite Cauchy--Schwarz applied to the singular
    values. The second is the defining application bound for the operator norm
    after Frobenius vectorization. Orthogonality makes the two rank-one
    projectors multiply to zero, so the squared Hilbert--Schmidt norm of their
    difference is the sum of their traces, namely $2$.
\end{proof}

\begin{theorem}[Convergence towards asymptotic states: upper bound]
    \label{thm:trace_norm_asymptotic_distance_upper}
    \lean{Matrix.traceNormAsymptoticDistance_pow_le_transferMatrix_l2_opNorm,
        Matrix.traceNormAsymptoticDistance_pow_le_hilbertSchmidtOperatorNorm}
    \leanok
    \uses{def:trace_norm_asymptotic_distance,
        def:peripheral_spectral_projections,
        def:mpu_transfer_matrix}
    Let $T:\MN{D}\to\MN{D}$ be positive and trace preserving, let $\rho$ be
    a density operator, and let $n>0$. Then
    \begin{align}
        \Delta_T(T^n(\rho))
        \leq \sqrt{\frac D2}\,
        \left\lVert \widehat T^{\,n}-\widehat T_\varphi^{\,n}\right\rVert_\infty\,
        \Delta_T(\rho).
        \tag{8.112}
    \end{align}
    Here the displayed transfer-matrix norm equals
    $\lVert T^n-T_\varphi^n\rVert_{2\to2}$, the Hilbert--Schmidt operator norm
    of the corresponding superoperator.
    The explicit $n>0$ records the paper's positive-natural convention: at
    $n=0$, both superoperator powers are the identity, so the displayed
    right-hand side vanishes. This is the upper assertion of~
    \cite[Chapter~8, Proposition ``Convergence towards asymptotic states'',
        Eqs.~(8.112), (8.114)--(8.116)]{Wolf2012Quantum}.
\end{theorem}

\begin{proof}\leanok
    \uses{lem:trace_norm_asymptotic_distance_numerator,
        lem:trace_norm_contraction_pointwise,
        lem:asymptotic_hilbert_schmidt_estimates,
        cor:peripheral_projection_channel,
        thm:transfer_matrix_hilbert_schmidt_operator_norm,
        lem:wolf_positive_power_nonperipheral_identity}
    Apply the numerator identity to $S=T^n-T_\varphi^n$. The peripheral
    projection is again positive and trace preserving, so Lemma~8.3 applies to
    the distinct density matrices $\rho$ and $T_\phi(\rho)$. Then use
    $\lVert\cdot\rVert_{\tr}\leq\sqrt D\lVert\cdot\rVert_2$ and the
    $\sqrt2$ Hilbert--Schmidt distance of orthogonal pure states. The constants
    simplify as $\frac12\sqrt D\sqrt2=\sqrt{D/2}$; if
    $\rho=T_\phi(\rho)$, both sides vanish directly. Finally, the
    positive-power identity identifies the transfer-matrix difference with
    the transfer matrix of $T^n-T_\varphi^n$, whose largest-singular-value norm
    equals the Hilbert--Schmidt operator norm
    $\lVert T^n-T_\varphi^n\rVert_{2\to2}$.
\end{proof}

\begin{lemma}[Four positive parts with Hilbert--Schmidt control]
    \label{lem:four_posSemidef_frobenius}
    \lean{Matrix.exists_four_posSemidef_frobenius_norm_le,
        Matrix.hilbertSchmidtOperatorNorm_le_of_forall}
    \leanok
    Every $X\in\MN{D}$ can be written $X=P_0-P_1+iP_2-iP_3$ with
    $P_j\succeq0$ and $\lVert P_j\rVert_2\leq\lVert X\rVert_2$.
\end{lemma}

\begin{proof}\leanok
    The Hermitian part $H=(X+X^\dagger)/2$ and the skew part
    $K=(X-X^\dagger)/(2i)$ have Hilbert--Schmidt norm at most
    $\lVert X\rVert_2$, and $X=H+iK$. For a Hermitian matrix with eigenvalues
    $\lambda_i$, the positive and negative parts are $f_\pm(H)$ with
    $f_+(x)=\max(x,0)$ and $f_-(x)=\max(-x,0)$; their squared
    Hilbert--Schmidt norms are $\sum_if_\pm(\lambda_i)^2\leq\sum_i\lambda_i^2$.
\end{proof}

\begin{theorem}[Convergence towards asymptotic states: lower bound]
    \label{thm:trace_norm_asymptotic_distance_lower}
    \lean{Matrix.hilbertSchmidtOperatorNorm_le_four_mul_of_posSemidef,
        Matrix.hilbertSchmidtOperatorNorm_pow_sub_le_of_forall_traceNormAsymptoticDistance_le,
        Matrix.hilbertSchmidtOperatorNorm_pow_sub_le_sSup_traceNormAsymptoticDistance,
        Matrix.transferMatrix_l2_opNorm_pow_sub_le_sSup_traceNormAsymptoticDistance}
    \leanok
    \uses{def:trace_norm_asymptotic_distance,
        def:peripheral_spectral_projections,
        def:mpu_transfer_matrix,
        lem:four_posSemidef_frobenius,
        thm:trace_norm_asymptotic_distance_upper}
    Let $S$ be a linear map on $\MN{D}$. If $\lVert S(P)\rVert_{\tr}\leq c$
    for every $P\succeq0$ with $\lVert P\rVert_2\leq1$, then
    \begin{align}
        \tfrac14\lVert S\rVert_{2\to2}\leq c.
        \tag{8.117}
    \end{align}
    Now let $T:\MN{D}\to\MN{D}$ be positive and trace preserving, let $D>0$,
    and let $n\in\N$. Then
    \begin{align}
        \sup_\rho\Delta_T(T^n(\rho))
        \geq\frac1{4\sqrt D}
        \left\lVert \widehat T^{\,n}-\widehat T_\varphi^{\,n}\right\rVert_\infty,
        \tag{8.113}
    \end{align}
    the supremum being over density operators. This is the converse assertion
    of~\cite[Chapter~8, Proposition ``Convergence towards asymptotic states'',
        Eqs.~(8.113), (8.117)]{Wolf2012Quantum}. The bound needs no positivity
    of $T$ once $c$ bounds every $\Delta_T(T^n(\rho))$; positivity and trace
    preservation make the supremum finite, by (8.112).
\end{theorem}

\begin{proof}\leanok
    \uses{lem:four_posSemidef_frobenius,
        thm:trace_norm_asymptotic_distance_upper,
        thm:transfer_matrix_hilbert_schmidt_operator_norm,
        lem:wolf_positive_power_nonperipheral_identity}
    For (8.117), decompose $X$ by Lemma~\ref{lem:four_posSemidef_frobenius}
    and use $\lVert\cdot\rVert_2\leq\lVert\cdot\rVert_{\tr}$ on each of the
    four terms. For (8.113) with $n>0$, take $S=T^n-T_\varphi^n$. Since
    $T_\varphi^n=T^nT_\phi$, one has $S(\rho)=T^n(\rho-T_\phi(\rho))$ and
    $\lVert S(\rho)\rVert_{\tr}=\Delta_T(T^n(\rho))$. A nonzero $P\succeq0$
    with $\lVert P\rVert_2\leq1$ has $\tr P\leq\sqrt D$, and $P/\tr P$ is a
    density operator, so $\lVert S(P)\rVert_{\tr}\leq\sqrt D\,c$ for every
    upper bound $c$ of the distances. For $n=0$ the superoperator difference
    vanishes. Finally, the transfer-matrix norm equals
    $\lVert T^n-T_\varphi^n\rVert_{2\to2}$.
\end{proof}

\begin{lemma}[Scaled positive-part trace estimate]
    \label{lem:scaled_posPart_trace_estimate}
    \lean{Matrix.re_trace_posPart_map_le_of_scaledTrace}
    \leanok
    \uses{def:positive_map_rectangular}
    Let $T:\MN{D}\to\MN{D'}$ be a positive linear map with
    $\tr[T(\rho)] = c\,\tr[\rho]$ for some nonnegative real $c$ and
    all $\rho\in\MN{D}$. Then for all Hermitian $H\in\MN{D}$,
    $\tr[(TH)^+]\leq c\cdot\tr[H^+]$.
    Generalizes Lemma~\ref{lem:positive_part_trace_estimate} from the
    trace-preserving case $c=1$.
\end{lemma}

\begin{proof}\leanok
    \uses{lem:posPart_support_projection_estimate}
    Lemma~\ref{lem:posPart_support_projection_estimate} gives
    $\tr[(TH)^+]\leq\tr[T(H^+)]$.  The trace-scaling hypothesis then
    yields $\tr[T(H^+)] = c\,\tr[H^+]$, giving the claimed bound.
\end{proof}

\begin{theorem}[Scaled trace-norm contractivity]
    \label{thm:scaled_trace_norm_contractivity}
    \lean{Matrix.traceNorm_map_le_of_positive_of_scaledTrace}
    \leanok
    \uses{def:trace_norm, def:positive_map_rectangular}
    Let $T:\MN{D}\to\MN{D'}$ be a positive linear map with
    $\tr[T(\rho)] = c\,\tr[\rho]$ for some nonnegative real $c$ and
    all $\rho\in\MN{D}$. Then for all Hermitian $H\in\MN{D}$,
    $\lVert T(H)\rVert_{\tr}\leq c\cdot\lVert H\rVert_{\tr}$.
    Generalizes Theorem~\ref{thm:trace_norm_contractivity} from the
    trace-preserving case $c=1$.
\end{theorem}

\begin{proof}\leanok
    \uses{lem:positive_map_hermitian, lem:scaled_posPart_trace_estimate, lem:trace_norm_jordan}
    Same argument as Theorem~\ref{thm:trace_norm_contractivity}: apply
    Lemma~\ref{lem:scaled_posPart_trace_estimate} to $H$ and $-H$, then
    sum using the Jordan trace-norm formula.
\end{proof}

\begin{theorem}[Quantum Doeblin contraction]
    \label{thm:quantum_doeblin_contraction}
    \lean{Matrix.traceNorm_map_sub_map_le_of_doeblin}
    \leanok
    \uses{def:trace_norm, def:positive_map_rectangular,
        def:trace_preserving_rectangular, def:density_matrices}
    Let $T,T':\MN{D}\to\MN{D'}$ be trace-preserving and
    Hermiticity-preserving linear maps with
    $T'(X)=\tr[X]\,Y$ for some fixed $Y\in\MN{D'}$.
    If $T-\varepsilon T'$ is positive for some $\varepsilon\geq0$, then
    for all density operators $\rho_1,\rho_2\in\MN{D}$,
    \begin{align}
        \lVert T(\rho_1)-T(\rho_2)\rVert_{\tr}
        \leq(1-\varepsilon)\,\lVert\rho_1-\rho_2\rVert_{\tr}.
        \label{eq:quantum_doeblin_contraction}
    \end{align}
    See \cite[Chapter~8, Theorem~8.17]{Wolf2012Quantum}.
\end{theorem}

\begin{proof}\leanok
    \uses{thm:scaled_trace_norm_contractivity}
    The map $S:=T-\varepsilon T'$ is positive by hypothesis.
    Since $T'(\rho_1-\rho_2)=\tr[\rho_1-\rho_2]\,Y=0$ (the densities
    have unit trace), $T(\rho_1)-T(\rho_2)=S(\rho_1-\rho_2)$. The trace
    of $S(\rho)$ is $(\tr[\rho])-\varepsilon\tr[\rho]=(1-\varepsilon)
        \tr[\rho]$, so $S$ satisfies the scaled-trace hypothesis of
    Theorem~\ref{thm:scaled_trace_norm_contractivity} with
    $c=1-\varepsilon$. Positivity of $S$ on $\rho_1$ forces
    $0\leq1-\varepsilon$ (by PSD trace nonnegativity), so the
    scaled-trace theorem applies and yields the claimed bound.
\end{proof}

\begin{definition}[Trace-prepare map]
    \label{def:trace_prepare_map}
    \lean{Matrix.tracePrepareMap}
    \lean{Matrix.tracePrepareMap_apply}
    \leanok
    For any matrix $Y\in\MN{D'}$, define
    $T'_Y:\MN{D}\to\MN{D'}$ by
    \begin{align}
        T'_Y(X) = \tr[X]\,Y.
        \label{eq:trace_prepare_map}
    \end{align}
\end{definition}

\begin{lemma}[Choi matrix of the trace-prepare map]
    \label{lem:choi_matrix_trace_prepare}
    \lean{Matrix.choiMatrix_tracePrepareMap}
    \leanok
    \uses{def:choi_matrix_rect, def:trace_prepare_map}
    For arbitrary $Y\in\MN{D'}$, the rectangular
    (output-factor-first) Choi matrix of $T'_Y$ is
    \begin{align}
        \tau(T'_Y) = \frac{1}{D}\,(Y\otimes\Id).
        \label{eq:choi_trace_prepare}
    \end{align}
    See \cite[Chapter~8, Eq.~(8.86)]{Wolf2012Quantum}.
\end{lemma}

\begin{proof}\leanok
    \uses{def:choi_matrix_rect, def:trace_prepare_map}
    The normalized omega slice is $D^{-1}E_{ij}$, where $E_{ij}$ is the
    matrix unit.  Since $\tr(E_{ij})=\delta_{ij}$, the trace of the slice
    is $D^{-1}\delta_{ij}$, and $T'_Y$ multiplies this scalar by $Y$.
    On the other hand, $(Y\otimes\Id)_{(a,i),(b,j)} = Y_{ab}\,\delta_{ij}$,
    so both sides equal $D^{-1}Y_{ab}\,\delta_{ij}$.
\end{proof}

\begin{theorem}[Choi-dominated Doeblin contraction]
    \label{thm:choi_dominated_doeblin}
    \lean{Matrix.traceNorm_map_sub_map_le_of_choi_domination}
    \leanok
    \uses{def:trace_norm, def:choi_matrix_rect,
        def:trace_preserving_rectangular, def:density_matrices}
    Let $T:\MN{D}\to\MN{D'}$ be a trace-preserving,
    Hermiticity-preserving linear map and $Y\in\MN{D'}$ Hermitian
    with $\tr[Y]=1$.  For $\varepsilon\geq0$, if
    the rectangular Choi matrix $\tau(T)$ satisfies
    \begin{align}
        \tau(T) \ge \frac{\varepsilon}{D}\,(Y\otimes\Id),
        \label{eq:choi_domination}
    \end{align}
    then for all density operators $\rho_1,\rho_2\in\MN{D}$,
    \begin{align}
        \lVert T(\rho_1)-T(\rho_2)\rVert_{\tr}
        \le (1-\varepsilon)\,\lVert\rho_1-\rho_2\rVert_{\tr}.
        \label{eq:choi_doeblin_contraction}
    \end{align}
    See \cite[Chapter~8, Eq.~(8.86)]{Wolf2012Quantum}.
\end{theorem}

\begin{proof}\leanok
    \uses{thm:quantum_doeblin_contraction, lem:choi_matrix_trace_prepare,
        thm:choi_rect_cp, def:trace_prepare_map, lem:is_kraus_cp_basic}
    Since $\tr[Y]=1$, $T'_Y$ is trace-preserving:
    $\tr[T'_Y(X)]=\tr[X]\tr[Y]=\tr[X]$.  Since $Y$ is Hermitian and
    $\tr[X]$ is real for Hermitian $X$, $T'_Y$ is Hermiticity-preserving.
    By Lemma~\ref{lem:choi_matrix_trace_prepare},
    $\tau(T'_Y)=(1/D)(Y\otimes\Id)$. Linearity of the Choi assignment gives
    $\tau(T-\varepsilon T'_Y)=\tau(T)-\varepsilon\tau(T'_Y)
        \ge0$, so $T-\varepsilon T'_Y$ is completely positive
    (the rectangular Choi CP equivalence, Theorem~\ref{thm:choi_rect_cp})
    and hence positive. Theorem~\ref{thm:quantum_doeblin_contraction} then
    yields the asserted bound.
\end{proof}

\begin{theorem}[Scalar Choi lower bound and traceless contraction]
    \label{thm:scalar_choi_traceless_contraction}
    \lean{Matrix.traceNorm_map_le_of_smul_one_le_choiMatrix}
    \leanok
    \uses{def:trace_norm, def:choi_matrix_rect}
    Let $D>0$ and let $T:\MN{D}\to\MN{D}$ be trace-preserving with Choi
    matrix $\tau$. If $c\,\Id\le\tau$ for a real $c$, then
    $0\le1-D^2c$ and every traceless Hermitian $H\in\MN{D}$ satisfies
    \begin{align}
        \lVert T(H)\rVert_{\tr}\le(1-D^2c)\,\lVert H\rVert_{\tr}.
    \end{align}
\end{theorem}

\begin{proof}\leanok
    \uses{lem:choi_matrix_trace_prepare, thm:choi_rect_cp,
        thm:scaled_trace_norm_contractivity}
    Take $Y=cD\,\Id$ in the trace-prepare map. Its Choi matrix is
    $D^{-1}(Y\otimes\Id)=c\,\Id$, so $S=T-T'_Y$ has Choi matrix
    $\tau-c\,\Id\ge0$ and is completely positive. It scales traces by
    $1-D^2c$, which is nonnegative because $S(\Id)\ge0$. On a traceless $H$
    the maps $S$ and $T$ agree, and a positive map with trace scale
    $1-D^2c$ contracts the trace norm of Hermitian matrices by that factor.
\end{proof}

\begin{corollary}[Channel spectrum versus Jamio{\l}kowski spectrum]
    \label{cor:channel_spectrum_jamiolkowski}
    \lean{Matrix.subperipheralModulus_le_of_smul_one_le_choiMatrix,
        Matrix.subperipheralModulus_le_one_sub_sq_mul_min_eigenvalue}
    \leanok
    \uses{def:wolf_subperipheral_modulus, def:choi_matrix_rect,
        def:is_kraus_cp}
    Let $T:\MN{D}\to\MN{D}$ be completely positive and trace-preserving
    with Jamio{\l}kowski state $\tau$. Let $\mu$ be the largest modulus of
    an eigenvalue of $T$ strictly inside the unit disc, with $\mu=0$ when
    there is none, and let $\lambda_{\min}(\tau)$ be the smallest eigenvalue
    of $\tau$. Then
    \begin{align}
        \mu\le1-D^2\lambda_{\min}(\tau).
    \end{align}
    More generally $\mu\le1-D^2c$ whenever $c\,\Id\le\tau$.
    This is \cite[Chapter~8, Equation~(8.87)]{Wolf2012Quantum}.
\end{corollary}

\begin{proof}\leanok
    \uses{thm:scalar_choi_traceless_contraction}
    For $D=0$ there is no eigenvector and $\mu=0\le1$. Otherwise write
    $q=1-D^2c\ge0$. Powers of $T$ preserve traceless Hermitian matrices,
    so $\lVert T^n(H)\rVert_{\tr}\le q^n\lVert H\rVert_{\tr}$ for each such
    $H$ by the preceding theorem. Let $T(X)=zX$ with $X\ne0$ and $|z|<1$.
    Trace preservation gives $z\tr X=\tr X$, hence $\tr X=0$. Complete
    positivity gives $T(X^\dagger)=T(X)^\dagger$, so the matrices
    $H=\tfrac12(X+X^\dagger)$ and $K=-\tfrac{i}{2}(X-X^\dagger)$ are traceless
    Hermitian with $X=H+iK$. Therefore
    \begin{align}
        |z|^n\lVert X\rVert_{\tr}=\lVert T^n(X)\rVert_{\tr}
        \le q^n\bigl(\lVert H\rVert_{\tr}+\lVert K\rVert_{\tr}\bigr)
    \end{align}
    for all $n$. If $|z|>q$, then $0\le q/|z|<1$, and dividing by $|z|^n$
    gives
    \begin{align}
        \lVert X\rVert_{\tr}\le\left(\frac{q}{|z|}\right)^n
            \bigl(\lVert H\rVert_{\tr}+\lVert K\rVert_{\tr}\bigr)
        \xrightarrow[n\to\infty]{}0,
    \end{align}
    so $\lVert X\rVert_{\tr}=0$, contradicting $X\ne0$. Hence $|z|\le q$ for every
    eigenvalue in the open unit disc, and $\mu\le q$. Finally
    $\lambda_{\min}(\tau)\,\Id\le\tau$. Wolf instead appeals to the asymptotic
    lower bound on $\lVert T^n(\rho_1-\rho_2)\rVert_{\tr}$; the eigenvector
    argument reaches the same conclusion directly.
\end{proof}

\section{Completely positive interpolation}

\begin{definition}[Condition number]
    \label{def:condition_number}
    \lean{Matrix.conditionNumber}
    \leanok
    For a square matrix $X$ the condition number is
    $\kappa(X)=\lVert X\rVert_\infty\lVert X^{-1}\rVert_\infty$, with the
    operator norm. See \cite[Chapter~8]{Wolf2012Quantum}.
\end{definition}

\begin{lemma}[Extremal eigenvalues of a positive definite matrix]
    \label{lem:pos_def_extremal_eigenvalues}
    \lean{Matrix.PosSemidef.exists_isUnitVector_mulVec_eq_norm_smul,
        Matrix.PosDef.inv_norm_inv_spec, Matrix.IsHermitian.le_norm_smul_one}
    \leanok
    Let $A\in\MN{d}$ be positive definite with $d>0$. Then
    $\lVert A^{-1}\rVert^{-1}\Id\le A\le\lVert A\rVert\Id$, the number
    $\lVert A^{-1}\rVert^{-1}$ is positive, and both $\lVert A\rVert$ and
    $\lVert A^{-1}\rVert^{-1}$ are eigenvalues of $A$ with unit eigenvectors.
\end{lemma}

\begin{proof}\leanok
    The norm of a positive element of a unital C$^*$-algebra lies in its
    spectrum and dominates it. Applied to $A$ this gives the upper bound and
    the top eigenvector. Applied to $A^{-1}$ it gives a unit $w$ with
    $A^{-1}w=\lVert A^{-1}\rVert w$, hence $Aw=\lVert A^{-1}\rVert^{-1}w$.
    For any eigenvector $v$ of $A$ with eigenvalue $\mu>0$, one has
    $\mu^{-1}=\langle v,A^{-1}v\rangle\le\lVert A^{-1}\rVert$, so every
    eigenvalue of $A$ is at least $\lVert A^{-1}\rVert^{-1}$, which is the
    lower bound.
\end{proof}

\begin{theorem}[Completely positive interpolation and the condition number]
    \label{thm:cp_unital_interpolation_condition_number}
    \lean{Matrix.exists_isKrausCP_unital_map_smul_iff_conditionNumber_le,
        Matrix.conditionNumber_le_of_positive_unital_map_smul}
    \leanok
    \uses{def:condition_number, def:is_kraus_cp}
    Let $P\in\MN{d}$ and $P'\in\MN{d'}$ be positive definite. There is a
    completely positive unital linear map $T:\MN{d}\to\MN{d'}$ with
    $T(P)=\lambda P'$ for some $\lambda>0$ if and only if
    \begin{align}
        \kappa(P)\ge\kappa(P').
    \end{align}
    The necessity of the inequality holds for every positive unital $T$.
    This is \cite[Chapter~8, Equations~(8.93)--(8.95)]{Wolf2012Quantum}.
\end{theorem}

\begin{proof}\leanok
    \uses{lem:pos_def_extremal_eigenvalues, lem:is_kraus_cp_basic,
        def:trace_prepare_map}
    Write $M=\lVert P\rVert$, $m=\lVert P^{-1}\rVert^{-1}$, and similarly
    $M',m'$ for $P'$, so that $\kappa(P)=M/m$ and $\kappa(P')=M'/m'$.
    Empty dimensions are checked directly: for $d'=0$ the zero map works,
    and for $d=0<d'$ neither side can hold.

    If $T$ is positive and unital with $T(P)=\lambda P'$, apply $T$ to the
    positive matrices $\Id-P/M$ and $P-m\Id$. This gives
    $\Id-(\lambda/M)P'\ge0$ and $\lambda P'-m\Id\ge0$. Evaluating at unit
    eigenvectors of $P'$ for $M'$ and $m'$ yields $\lambda M'\le M$ and
    $m\le\lambda m'$, so $M'/m'\le M/m$.

    Conversely let $M'/m'\le M/m$ and $\lambda=M/M'$. If $m<M$, let
    $\psi_1,\psi_2$ be unit eigenvectors of $P$ for $m$ and $M$ and put
    \begin{align}
        T(X)=\langle\psi_1|X|\psi_1\rangle\,\frac{M\Id-\lambda P'}{M-m}
            +\langle\psi_2|X|\psi_2\rangle\,\frac{\lambda P'-m\Id}{M-m}.
    \end{align}
    The first coefficient matrix is $\lambda(M'\Id-P')/(M-m)\ge0$, and the
    second is positive because $\lambda m'\ge m$ is the hypothesis. Each
    summand is a rank-one compression followed by preparation of a positive
    semidefinite matrix, so $T$ is completely positive, and direct
    evaluation gives $T(\Id)=\Id$ and $T(P)=\lambda P'$. This is Wolf's map,
    since $E_1(X)=\Id-X/M$ and $E_2(X)=X-m\Id$ give
    $\langle\psi_1|E_1(P)|\psi_1\rangle=(M-m)/M$ and
    $\langle\psi_2|E_2(P)|\psi_2\rangle=M-m$. If $m=M$ then $P=M\Id$, the
    hypothesis forces $m'=M'$ and $P'=M'\Id$, and
    $X\mapsto\langle\psi|X|\psi\rangle\Id$ works for any unit $\psi$.
\end{proof}

\section{Projective pinching}

\begin{theorem}[Projective pinching inequality]
    \label{thm:projective_povm_pinching}
    \lean{projectivePOVM_pinching}
    \leanok
    \uses{def:density_matrices, def:orthogonal_projection_channel}
    Let $\rho\in\MN{d}$ be a density matrix and let
    $P_1,\ldots,P_k\in\MN{d}$ be orthogonal projections satisfying
    $\sum_{i=1}^k P_i=\Id$.  Define
    \begin{align}
        T(X) & = \sum_{i=1}^k P_iXP_i .
    \end{align}
    Then
    \begin{align}
        \rho & \le kT(\rho).
        \label{eq:projective_pinching}
    \end{align}
    This is \cite[Chapter~8, Eq.~(8.56)]{Wolf2012Quantum}.
\end{theorem}

\begin{proof}\leanok
    For $A_{ij}=P_i-P_j$, positivity of $\rho$ gives
    $A_{ij}\rho A_{ij}^{\dagger}\ge 0$.  Since the projections are
    Hermitian and sum to the identity, expansion and collection of the four
    double sums gives
    \begin{align}
        \sum_{i,j=1}^k A_{ij}\rho A_{ij}^{\dagger}
         & = 2\bigl(kT(\rho)-\rho\bigr).
    \end{align}
    Thus $kT(\rho)-\rho$ is positive semidefinite.
\end{proof}

\section{Cone ratios and norm comparisons}

\begin{definition}[Lower cone ratio]
    \label{def:lower_cone_ratio}
    \lean{ProperCone.lowerRatio, ProperCone.NormMonotone}
    \leanok
    Let $C$ be a closed convex cone in a real normed space, and suppose
    $\|x\|\le\|x+y\|$ for $x,y\in C$. For $a,b\in C$ with $b\ne0$, put
    $\inf(a/b)=\sup\{t\ge0:a-tb\in C\}$.
    Here the source notation is defined by the order description in
    \cite[Chapter~8, Equation~(8.61)]{Wolf2012Quantum}.
\end{definition}

\begin{lemma}[Attainment and bounds for cone ratios]
    \label{lem:lower_cone_ratio_bounds}
    \lean{ProperCone.lowerRatio_bddAbove, ProperCone.lowerRatio_mem,
        ProperCone.le_lowerRatio, ProperCone.lowerRatio_le_norm_div,
        ProperCone.lowerRatio_mul_le_one}
    \leanok
    \uses{def:lower_cone_ratio}
    Let $b\in C$ with $b\ne0$, and let $a$ be any vector in the ambient space.
    The admissible set $S=\{t\ge0:a-tb\in C\}$ is bounded above by
    $\|a\|/\|b\|$. Every $t\in S$ satisfies $t\le\inf(a/b)$.
    If additionally $a\in C$, then $S$ contains zero, its supremum is attained,
    and $0\le\inf(a/b)\le\|a\|/\|b\|$.
    If also $a\ne0$, then $\inf(a/b)\inf(b/a)\le1$.
\end{lemma}
\begin{proof}\leanok
    Monotonicity gives $t\|b\|\le\|a\|$ for each admissible $t$.
    The supremum bounds every admissible scale. When $a\in C$, zero is
    admissible, and closedness gives attainment. When also $a\ne0$,
    multiplying the norm bounds in both directions gives the product bound.
\end{proof}

\begin{theorem}[Norm comparison from cone ratios]
    \label{thm:cone_ratio_norm_bound}
    \lean{ProperCone.norm_sub_le_lowerRatio, ProperCone.norm_sub_le_of_scales,
        ProperCone.norm_sub_le_lowerRatio_average, ProperCone.norm_sub_le_lowerRatio_sqrt,
        ProperCone.norm_sub_le_lowerRatio_log}
    \leanok
    \uses{lem:lower_cone_ratio_bounds}
    For nonzero $a,b\in C$, write $\nu=\inf(b/a)$ and $\mu=\inf(a/b)$. Then
    \begin{align}
        \|a-b\|\le |1-\nu|\|a\|+|1-\mu|\|b\|
            +\min\{|1-\nu|\|a\|,|1-\mu|\|b\|\}.
    \end{align}
    If $\mu,\nu\le1$, then
    \begin{align}
        \|a-b\|&\le\frac32(\|a\|+\|b\|-\nu\|a\|-\mu\|b\|)\\
        &\le3\left(\frac{\|a\|+\|b\|}{2}
             -\sqrt{\nu\|a\|\mu\|b\|}\right).
    \end{align}
    If also $\mu\nu>0$, the square-root term equals
    $\sqrt{\|a\|\|b\|}\exp(\log(\mu\nu)/2)$.
    These are \cite[Chapter~8, Equations~(8.63), (8.68)--(8.69)]{Wolf2012Quantum},
    with the last expression written directly in terms of the lower ratios.
\end{theorem}
\begin{proof}\leanok
    Put $p=b-\nu a$ and $q=a-\mu b$. Both belong to $C$, and
    $p+q=(1-\nu)a+(1-\mu)b$. The triangle inequality and norm monotonicity give
    $\|a-b\|\le2|1-\nu|\|a\|+|1-\mu|\|b\|$.
    Interchanging $a,b$ and taking the minimum gives the first bound.
    When $\mu,\nu\le1$, averaging the two asymmetric estimates gives
    \begin{align}
        \|a-b\|&\le\frac32\bigl(\|a\|+\|b\|-x-y\bigr),
        \label{eq:cone_norm_average}
    \end{align}
    where $x=\nu\|a\|\ge0$ and $y=\mu\|b\|\ge0$. Since
    \begin{align}
        x+y-2\sqrt{xy}&=(\sqrt{x}-\sqrt{y})^2\ge0,
        \label{eq:cone_norm_amgm}
    \end{align}
    substituting (\ref{eq:cone_norm_amgm}) into (\ref{eq:cone_norm_average}) yields
    \begin{align}
        \|a-b\|&\le3\left(\frac{\|a\|+\|b\|}{2}
            -\sqrt{\nu\|a\|\mu\|b\|}\right).
    \end{align}
\end{proof}

\begin{theorem}[Cone ratios under positive maps]
    \label{thm:cone_ratio_positive_map}
    \lean{ProperCone.lowerRatio_le_map, ProperCone.neg_log_lowerRatio_map_le}
    \leanok
    \uses{lem:lower_cone_ratio_bounds}
    Let $K$ be a closed convex cone in another real normed space whose norm
    is monotone on $K$, and let $T$ be linear with $T(C)\subseteq K$.
    For $a,b\in C$ with $b\ne0$ and $Tb\ne0$,
    $\inf(a/b)\le \inf(Ta/Tb)$. If additionally $a\ne0$, $Ta\ne0$,
    and both input ratios are positive, then
    \begin{align}
        -\log\bigl(\inf(Ta/Tb)\inf(Tb/Ta)\bigr)
          \le -\log\bigl(\inf(a/b)\inf(b/a)\bigr).
    \end{align}
\end{theorem}
\begin{proof}\leanok
    Apply $T$ to $a-\inf(a/b)b\in C$ and use the supremum property in $K$.
    Multiply the inequalities in both directions and use monotonicity of
    the logarithm.
\end{proof}

\begin{definition}[Expanding diagonal map on the quadrant]
    \label{def:expanding_quadrant_map}
    \lean{ProjectiveDiameterCounterexample.quadrant,
        ProjectiveDiameterCounterexample.diagonalMap,
        ProjectiveDiameterCounterexample.mem_quadrant}
    \leanok
    Let $C=\mathbb R_+^2$ and $T(x,y)=(x,2y)$.
\end{definition}

\begin{theorem}[Boundary fixed point and an expanding eigenvector]
    \label{thm:boundary_fixed_point_expansion}
    \lean{ProjectiveDiameterCounterexample.boundary_fixed_point_counterexample,
        ProjectiveDiameterCounterexample.diagonalMap_positive,
        ProjectiveDiameterCounterexample.eigenvalue_gt_tanh,
        ProjectiveDiameterCounterexample.lowerRatio_first,
        ProjectiveDiameterCounterexample.lowerRatio_second,
        ProjectiveDiameterCounterexample.output_log_distance_unbounded}
    \leanok
    \uses{def:expanding_quadrant_map, def:lower_cone_ratio}
    The positive map $T$ fixes the nonzero boundary vector $(1,0)$ and has
    eigenvalue $2$ on $(0,1)$, which is not proportional to the fixed vector.
    The eigenvalue exceeds every coefficient at most one, including every
    $\tanh(D/4)$ for real $D$. Moreover, the logarithmic distances between
    images of interior vectors are unbounded.
    This obstructs the unrestricted extension proposed after
    \cite[Chapter~8, Equation~(8.91)]{Wolf2012Quantum} if the convention
    $\tanh(+\infty)=1$ is adopted. Without such a convention, the expression
    is undefined for this map. The example does not refute the theorem with
    finite diameter or an interior fixed point.
\end{theorem}
\begin{proof}\leanok
    Positivity and the two eigenvector identities follow by coordinates.
    For $t\ge1$, the images of $(1,t/2)$ and $(1,1/2)$ are $(1,t)$ and $(1,1)$.
    Their lower cone ratios are $1$ and $1/t$. Hence their logarithmic
    distance is $\log t$, which exceeds $R$ when
    $t=\exp(\max\{R,0\}+1)$.
\end{proof}
